%----------------------------------------------------------------------------
\chapter*{Összefoglaló}\addcontentsline{toc}{chapter}{Összefoglaló}
%----------------------------------------------------------------------------

A felgyorsult digitalizáció rengeteg újítást hoz a modern ember életében, beszivárog a mindennapjainka, legyen szó munkahelyről, oktatásról vagy épp egyéb területről. 

A Project Raccoon a digitalizáció által nyújtott lehetőségeket kihasználva továbbfejleszti az ügyintézési folyamatokat, segítve ezzel a modern ember mindennapjait. A platform rugalmas és felhasználóbarát felülete révén könnyen kezelhető, és a felhasználók számára egyszerűbbé teszi az adminisztratív feladatokat.

Ez a rendszer számos előnyt kínál a felhasználók számára, akik így időt és erőforrásokat takaríthatnak meg azáltal, hogy a hagyományos papíralapú folyamatokat digitális formába helyezik át. Segítségével az emberek könnyedén intézhetik ügyeiket otthonról, anélkül, hogy személyesen kellene felkeresniük különböző hivatalokat vagy intézményeket. Az online felületen elérhetővé válnak az igazolások, űrlapok, dokumentumok, amelyeket könnyedén kitölthetnek. Ezenkívül a platform lehetővé teszi a dokumentumok elektronikus aláírását, amely hitelesített módon helyettesíti a hagyományos, papíralapú aláírást. Ezáltal a felhasználók biztonságosan és megbízhatóan intézhetik ügyeiket az online térben.

A mérésekből bebizonyosodott, hogy mind a hagyományos szerver alapú kitelepítéseknek, mind a \textit{serverless} technológiákkal kitelepített környezeteknek még mindig van létjogosultságuk. Attól függően, hogy mekkora projektet készítünk, és milyen gyakran fogják használni a felhasználók, kell döntenünk a hasonló telepítési kérdésekről.

\hiddensubsection{GitHub projekt linkje}

A projekt forráskódja elérhetően a következő linken: \url{https://github.com/darktohka/project-raccoon}.